\documentclass[11pt]{article}
\usepackage{amsmath,amssymb,amsfonts}
\usepackage{geometry}
\geometry{margin=1in}
\usepackage{graphicx}
\usepackage{booktabs}
\usepackage{natbib}
\bibliographystyle{plainnat}

\newcommand{\bxi}{\boldsymbol{\xi}}
\newcommand{\bx}{\mathbf{x}}
\newcommand{\avg}[1]{\langle #1 \rangle}

\title{Metastability of Retrieval in LSR Dense Associative Memory}
\author{}
\date{}

\begin{document}
\maketitle

\section{Introduction}

Associative memory networks store patterns so that the correct memory can be retrieved from a noisy or partial cue.  In the classical Hopfield network~\citep{Hopfield1982}, storage capacity scales linearly with the number of neurons~\citep{amit1987statistical}.  Dense Associative Memory (DAM) models replace pairwise interactions with highly nonlinear energy functions, achieving exponential storage: $M = e^{\alpha N}$ patterns for $N$ continuous neurons on the unit sphere, where $\alpha = \ln M/N$ is the logarithmic load per neuron~\citep{demircigil2017model,ramsauer2021hopfield}.

Two kernel-based energy functions have been studied in this setting.  The log-sum-exp (LSE) energy, based on a Gaussian similarity kernel, was analyzed by~\citet{ramsauer2021hopfield}, who proved a zero-temperature capacity $\alpha_c(0) = 0.5$.  More recently, \citet{hoover2025dense} introduced the log-sum-ReLU (LSR) energy for binary neurons, based on the compactly supported Epanechnikov kernel.  For continuous neurons on the $N$-sphere ($|\bx|^2 = N$), the LSR energy takes the form
\begin{equation}
E(\bx) = -\frac{N}{b}\ln\left(\sum_{\mu=1}^{M}\max\!\Big(0,\;1 - b + b\,\phi_\mu\Big)\right),
\qquad \phi_\mu \equiv \frac{\bxi^\mu \cdot \bx}{N}\,,
\label{eq:lsr}
\end{equation}
where $b > 1$ is the rescaled inverse variance controlling the kernel sharpness, $\bxi^\mu$ are the stored patterns on $S^{N-1}(\!\sqrt{N})$, and $\phi_\mu$ is the alignment of the state $\bx$ with pattern~$\mu$.  The compact support introduces a hard wall at $\varphi_c = (b-1)/b$: pattern~$\mu$ contributes to the energy only when $\phi_\mu > \varphi_c$.  This is qualitatively different from the LSE kernel, where all patterns contribute everywhere on the sphere.

The thermodynamic phase boundaries $\alpha_c(T)$ for both kernels on the $N$-sphere were derived by~\citet{icml2026theory} using the replica method, and validated by Monte Carlo (MC) simulations in~\citet{iclr2026mc}.  A key result for the LSR kernel is the existence of a support threshold
\begin{equation}
\alpha_{\mathrm{th}} = \frac{\varphi_c^2}{2} = \frac{1}{2}\!\left(\frac{b-1}{b}\right)^{\!2},
\label{eq:alpha_th}
\end{equation}
below which no spurious pattern falls within the support and retrieval is perfect at any temperature~\citep{icml2026theory}.  For $b = 2+\sqrt{2}$ (the value used throughout this paper): $\varphi_c \approx 0.707$ and $\alpha_{\mathrm{th}} = 0.25$.  This sub-threshold regime has no analog in the LSE model, where a critical line exists at all loads $\alpha > 0$.

The derivation of~\eqref{eq:alpha_th} relies on approximating the alignment distribution on $S^{N-1}$ by a Gaussian $f(\phi) \propto e^{-N\phi^2/2}$, as in~\citet{ramsauer2021hopfield}.  The exact density on the sphere,
\begin{equation}
f(\phi) = C_N\,(1 - \phi^2)^{(N-3)/2}\,,\qquad
C_N = \frac{\Gamma(N/2)}{\Gamma\!\big((N\!-\!1)/2\big)\,\sqrt{\pi}}\,,
\label{eq:exact_density}
\end{equation}
has the same mean ($0$) and variance ($1/N$) but lighter tails: it vanishes at $|\phi| = 1$, whereas the Gaussian assigns nonzero weight there.  Since the LSR noise floor depends on whether any pattern exceeds $\varphi_c$, the lighter tails of~\eqref{eq:exact_density} predict fewer patterns above threshold and shift $\alpha_{\mathrm{th}}$ to a higher value ($\approx 0.347$ for $b = 2+\sqrt{2}$).

\begin{small}
\noindent\textit{Remark (to be removed).}  Let $\varphi_{1,\max}$ denote the maximum of the overlap $\bxi^1\!\cdot\!\bxi^\mu/N$ over all spurious patterns $\mu \neq 1$; it is defined by the condition that on average one pattern reaches this overlap.  The Gaussian gives $\varphi_{1,\max} = \sqrt{2\alpha}$; the exact density~\eqref{eq:exact_density} gives $\varphi_{1,\max}^{\mathrm{exact}} = \sqrt{1-e^{-2\alpha}}$.  Setting $\varphi_{1,\max}^{\mathrm{exact}} = \varphi_c$: $\alpha_{\mathrm{th}}^{\mathrm{exact}} = -\tfrac{1}{2}\ln(1-\varphi_c^2) = (\ln 2)/2 \approx 0.347$.  Finite-$N$ corrections shift this further right: $\alpha_{\mathrm{th}}(N) = \alpha_{\mathrm{th}}^{\mathrm{exact}} + [\tfrac{1}{2}\ln N + \mathrm{const}]/N \approx 0.39$ at $N = 50$.
\end{small}

\noindent However, the MC simulations of~\citet{iclr2026mc} place the LSR retrieval boundary much closer to $\alpha \approx 0.25$ than to~$0.347$.  The Gaussian approximation, despite its unphysical tails, describes the onset of interference more accurately than the exact density~\eqref{eq:exact_density}.  The reason for this is not understood; we note only that finite-$N$ corrections to the exact prediction push $\alpha_{\mathrm{th}}$ further from the MC value, not closer (see the Remark above).

In this paper we set aside the question of the precise location of $\alpha_{\mathrm{th}}$ and instead examine what happens \emph{inside} the retrieval region $\alpha < \alpha_{\mathrm{th}}$, where the thermodynamic theory of~\citet{icml2026theory} predicts perfect retrieval at any temperature.  We reproduce and extend the MC simulations of~\citet{iclr2026mc} for the LSR kernel, recording per-trial observables (two replicas $\times$ up to 256 independent disorder samples per $(\alpha, T)$ point) at finite system sizes $N \approx 20$--$50$.

The per-trial overlap distributions reveal that the retrieval region is richer than the thermodynamic prediction suggests.  Figure~\ref{fig:cdf} shows the empirical CDF of the time-averaged alignment $\varphi$ with the target pattern, pooled over disorder realizations and both replicas.  Rather than a single peak at the retrieval overlap $\phi_{\mathrm{eq}}(T)$, the distribution is multimodal: distinct concentrations appear at $\varphi \approx 0.90$, $0.83$, $0.48$, and near~$0$.  Each peak corresponds to a different (meta)stable state that the system can reach within the allocated MC time.  The weight of the non-retrieval peaks grows with~$\alpha$ and with the number of MC steps, suggesting that the retrieval state is metastable rather than stable at these finite system sizes.

The goal of this paper is to explain these peaks---their positions, relative weights, and dependence on $\alpha$, $T$, and $N$---and to identify the mechanism by which the system escapes from the retrieval basin.  The analysis draws comparisons with the classical Hopfield model~\citep{amit1987statistical}, where mixture states and spin-glass phases are known to coexist with retrieval.

\begin{figure}[t]
\centering
\includegraphics[width=\textwidth]{panels_v8/cdf_phi.pdf}
\caption{Empirical CDF of the per-trial time-averaged alignment $\varphi$ with the target pattern (both replicas pooled).  Green: $(\alpha, T)$ points in the retrieval region; orange: intermediate; blue: non-retrieval.  Dashed gray line: $\varphi_c \approx 0.707$.  The step structure reflects concentration of trials at discrete overlap values.}
\label{fig:cdf}
\end{figure}

\paragraph{Outline.}  Section~\ref{sec:theory} develops the theory: support caps and their intersections, the two-pattern centroid and its energy, multi-pattern centroids as a function of~$b$, the zero-energy escape path, the Kramers entropic barrier, the two-pattern thermal equilibrium, and centroid hopping.  Section~\ref{sec:mc} presents the Monte Carlo simulations, the observed overlap distributions, and their quantitative comparison with the theory.  Section~\ref{sec:discussion} summarizes and discusses the implications for retrieval robustness.

%%% ─────────────────────────────────────────────────────────────────
\section{Theory of the centroid escape}
\label{sec:theory}

We use the following results from~\citet{icml2026theory} without re-derivation: the internal energy density $u_{\mathrm{LSR}}(\phi) = -b^{-1}\ln[1-b(1-\phi)]$ (ibid., Eq.~14), the geometric entropy $s(\phi) = \tfrac{1}{2}\ln(1-\phi^2)$ (ibid., Eq.~19), and the single-pattern self-consistency equation for the equilibrium alignment $\phi_{\mathrm{eq}}(T)$ (ibid., Eq.~37).  At $b = 2+\sqrt{2}$ and $T = 0.325$: $\phi_{\mathrm{eq}} \approx 0.90$.

\subsection{Support caps on the sphere}
\label{ssec:caps}

Each pattern $\bxi^\mu$ defines a \emph{support cap} on the sphere:
\begin{equation}
S_\mu = \big\{\bx \in S^{N-1}(\!\sqrt{N}) : \phi_\mu(\bx) > \varphi_c\big\}\,,
\label{eq:cap}
\end{equation}
the set of states for which pattern~$\mu$ contributes to the LSR energy~\eqref{eq:lsr}.  The cap is centered at $\bxi^\mu$ with angular half-width $\theta_c = \arccos(\varphi_c)$.  Increasing $b$ raises $\varphi_c$ and shrinks the caps.

The \emph{inter-pattern overlap} $\varphi_{1\mu} = \bxi^1\!\cdot\!\bxi^\mu/N$ measures the cosine of the angle between patterns $\bxi^1$ and $\bxi^\mu$.  Two caps $S_1$ and $S_\mu$ intersect if and only if this angle is less than $2\theta_c$, i.e.,
\begin{equation}
\varphi_{1\mu} > \cos(2\theta_c) = 2\varphi_c^2 - 1\,.
\label{eq:cap_intersection}
\end{equation}
The threshold $2\varphi_c^2 - 1$ depends on $b$ through $\varphi_c = (b-1)/b$: it equals~$0$ at $b = 2+\sqrt{2}$, and is positive for larger~$b$.

For random patterns on $S^{N-1}(\!\sqrt{N})$, $\varphi_{1\mu}$ is distributed according to the exact density~\eqref{eq:exact_density}, with mean~$0$ and variance~$1/N$.  With $M = e^{\alpha N}$ patterns, the expected number having overlap exceeding any threshold $\varphi^*$ is
\begin{equation}
n(\varphi^*) = M\,\Pr(\varphi_{1\mu} > \varphi^*)\,.
\label{eq:n_count}
\end{equation}
The maximum inter-pattern overlap $\varphi_{1,\max}$, defined by $n(\varphi_{1,\max}) = 1$, is $\sqrt{2\alpha}$ in the Gaussian approximation and $\sqrt{1-e^{-2\alpha}}$ using the exact density (see Section~1).

\subsection{The two-pattern centroid}
\label{ssec:centroid}

Consider the target pattern $\bxi^1$ and a spurious pattern $\bxi^\mu$ with inter-pattern overlap $\varphi_{1\mu} > 0$.  We work in the two-dimensional subspace spanned by the orthonormalized pair
\begin{equation}
\mathbf{e}_1 = \frac{\bxi^1}{\sqrt{N}}\,, \qquad
\mathbf{e}_2 = \frac{\bxi^\mu/\sqrt{N} - \varphi_{1\mu}\,\mathbf{e}_1}{\sqrt{1-\varphi_{1\mu}^2}}\,.
\label{eq:basis}
\end{equation}
Projecting the state $\bx/\sqrt{N} = a\,\mathbf{e}_1 + v\,\mathbf{e}_2 + \mathbf{c}_\perp$, with $a^2 + v^2 + |\mathbf{c}_\perp|^2 = 1$, the alignments with the two patterns are
\begin{equation}
\phi_1 = a\,, \qquad \phi_\mu = a\,\varphi_{1\mu} + v\sqrt{1-\varphi_{1\mu}^2}\,.
\label{eq:overlaps}
\end{equation}

The support conditions $\phi_1 > \varphi_c$ and $\phi_\mu > \varphi_c$ define two half-planes in the $(a,v)$ disk.  Their intersection---the region where both patterns are simultaneously in support---is non-empty whenever $\varphi_{1\mu} > 2\varphi_c^2 - 1$ (Eq.~\eqref{eq:cap_intersection}).

Within this intersection, the LSR energy~\eqref{eq:lsr} is
\begin{equation}
\frac{E}{N} = -\frac{1}{b}\ln\!\big[(1-b+b\,\phi_1) + (1-b+b\,\phi_\mu)\big]\,.
\label{eq:E_two}
\end{equation}
At zero temperature, the energy is minimized at the symmetric point $\phi_1 = \phi_\mu$, which from~\eqref{eq:overlaps} gives the \emph{geometric centroid}:
\begin{equation}
\varphi_{\mathrm{cen}}^{(0)} = \sqrt{\frac{1+\varphi_{1\mu}}{2}}\,.
\label{eq:centroid}
\end{equation}

\paragraph{Centroid energy vs retrieval energy.}
The energy at the geometric centroid (two patterns in support) compared with the retrieval state (one pattern):
\begin{equation}
\frac{E_{\mathrm{ret}}}{N} = u_{\mathrm{LSR}}(\phi_{\mathrm{eq}})\,,\qquad
\frac{E_{\mathrm{cen}}^{(0)}}{N} = -\frac{1}{b}\ln\!\Big(2\big[1-b+b\,\varphi_{\mathrm{cen}}^{(0)}\big]\Big)\,.
\label{eq:E_compare}
\end{equation}
At $\alpha = 0.20$, the exact density~\eqref{eq:exact_density} gives $\varphi_{1,\max} = \sqrt{1-e^{-0.4}} \approx 0.57$.  With $\phi_{\mathrm{eq}} = 0.90$: $E_{\mathrm{ret}}/N = 0.12$, $E_{\mathrm{cen}}^{(0)}/N = -0.01$.  The centroid has \emph{lower energy} than the retrieval state.

\paragraph{$k$-pattern centroids.}
If $k$ patterns have a common pairwise overlap $\varphi_{1\mu}$, the symmetric $k$-pattern centroid has geometric overlap $\varphi_k^{(0)} = \sqrt{\varphi_{1\mu} + (1-\varphi_{1\mu})/k}$ and energy $E_k/N = -b^{-1}\ln[k(1-b+b\,\varphi_k^{(0)})]$.  As $k \to \infty$, $\varphi_k^{(0)} \to \sqrt{\varphi_{1\mu}}$.  Whether adding patterns lowers the energy depends on whether this limiting overlap stays in support:
\begin{equation}
\sqrt{\varphi_{1\mu}} > \varphi_c \qquad\Longleftrightarrow\qquad \varphi_{1\mu} > \varphi_c^2\,.
\label{eq:cluster}
\end{equation}
For $b = 2+\sqrt{2}$: $\varphi_c^2 = 0.50$.  Since $\varphi_{1,\max} < \varphi_c$ in the sub-threshold regime $\alpha < \alpha_{\mathrm{th}}$ (by definition), and $\varphi_c^2 < \varphi_c$, multi-pattern centroids may or may not be favored depending on whether $\varphi_{1,\max}$ exceeds $\varphi_c^2$.  At $b = 2+\sqrt{2}$ and $\alpha = 0.20$: $\varphi_{1,\max} \approx 0.57 > 0.50$, so tight clusters of patterns with all pairwise overlaps $> 0.50$ would form deeper centroids.  However, such clusters require all $\binom{k}{2}$ pairwise overlaps to simultaneously exceed $\varphi_c^2$, which is exponentially rare at finite $N$.  In practice, the system hops between \emph{pairs}, not clusters.

For larger $b$, the cluster threshold $\varphi_c^2$ rises (e.g.\ $0.64$ at $b = 5$, $0.81$ at $b = 10$), making multi-pattern centroids even rarer.

\subsection{The escape path and the Kramers barrier}
\label{ssec:escape}

\paragraph{Stage 1: perpendicular rotation ($\Delta E = 0$).}
At the retrieval state, $a = \phi_{\mathrm{eq}}$ and $v = 0$.  The alignment with $\bxi^\mu$ is $\phi_\mu = \phi_{\mathrm{eq}}\,\varphi_{1\mu} < \varphi_c$: only the target is in support.  The system can increase $v$ (rotating $\bx_\perp$ toward $\bxi^\mu_\perp$) while keeping $a = \phi_1$ fixed.  Since only $\bxi^1$ is in support and $\phi_1$ does not change, the energy is exactly constant along this rotation.

When $v$ reaches the entry value
\begin{equation}
v_{\mathrm{entry}} = \frac{\varphi_c - \phi_{\mathrm{eq}}\,\varphi_{1\mu}}{\sqrt{1-\varphi_{1\mu}^2}}\,,
\label{eq:ventry}
\end{equation}
the spurious pattern enters support ($\phi_\mu = \varphi_c$).

\paragraph{Escape-path constraint.}
The rotation is possible only if $a^2 + v_{\mathrm{entry}}^2 \leq 1$ (the state must remain on the sphere).  This requires
\begin{equation}
\phi_{\mathrm{eq}}^2 + \frac{(\varphi_c - \phi_{\mathrm{eq}}\,\varphi_{1\mu})^2}{1-\varphi_{1\mu}^2} \;\leq\; 1\,,
\label{eq:sphere_constraint}
\end{equation}
which defines a minimum inter-pattern overlap for the escape:
\begin{equation}
\varphi_{\min}(\phi_{\mathrm{eq}}) = \varphi_c\,\phi_{\mathrm{eq}} - \sqrt{(1-\varphi_c^2)(1-\phi_{\mathrm{eq}}^2)}\,.
\label{eq:phi_min}
\end{equation}
For $b = 2+\sqrt{2}$ and $\phi_{\mathrm{eq}} = 0.90$: $\varphi_{\min} = 0.328$.  We denote by $K = n(\varphi_{\min})$ the number of escape-enabling patterns, i.e.\ the expected number of spurious patterns with $\varphi_{1\mu} > \varphi_{\min}$ (Eq.~\eqref{eq:n_count}).  Table~\ref{tab:n_overlap} gives $K$ at representative $\alpha$ values, computed by numerical integration of the exact density~\eqref{eq:exact_density}.

\begin{table}[h]
\centering
\caption{Expected number of patterns with $\varphi_{1\mu} > \varphi_{\min} = 0.328$ ($\phi_{\mathrm{eq}} = 0.90$, $b = 2+\sqrt{2}$), computed from the exact density~\eqref{eq:exact_density}.}
\label{tab:n_overlap}
\begin{tabular}{ccrrr}
\toprule
$\alpha$ & $N$ & $M$ & $n(\varphi_{1\mu} > 0.328)$ \\
\midrule
0.05 & 198 & 20\,000 & $0.02$ \\
0.10 &  99 & 20\,000 & 8.6 \\
0.20 &  50 & 20\,000 & 188 \\
0.30 &  36 & 43\,600 & 842 \\
\bottomrule
\end{tabular}
\end{table}

\paragraph{Stage 2: downhill slide ($\Delta E < 0$).}
Once both patterns are in support, the sum inside the logarithm in~\eqref{eq:lsr} increases and the energy decreases.  The system slides to the two-pattern centroid~\eqref{eq:centroid}.  The energy barrier along the full escape path is \emph{zero}: Stage~1 has $\Delta E = 0$ and Stage~2 has $\Delta E < 0$.

\paragraph{The entropic barrier.}
What prevents instantaneous escape is entropy, not energy.  The rotation in Stage~1 must align the $(N\!-\!1)$-dimensional perpendicular component $\bx_\perp$ with a specific direction ($\bxi^\mu_\perp$).  The perpendicular component lives on a sphere of radius $\sqrt{N(1-a^2)}$; increasing $v$ from $0$ to $v_{\mathrm{entry}}$ reduces the available perpendicular volume.  The free energy cost of this alignment is
\begin{equation}
\frac{\Delta F}{T} = \frac{N-2}{2}\left[-\ln\!\left(1 - \frac{v_{\mathrm{entry}}^2}{1-\phi_{\mathrm{eq}}^2}\right)\right] \;\approx\; \frac{N}{2}\,\frac{v_{\mathrm{entry}}^2}{1-\phi_{\mathrm{eq}}^2}\,.
\label{eq:kramers_barrier}
\end{equation}
This is the barrier of a Kramers escape problem.  The variable $v$ diffuses in an entropic harmonic potential (from the sphere constraint), and the escape requires a thermal fluctuation to $v = v_{\mathrm{entry}}$.  The escape time for a single pattern scales as
\begin{equation}
\tau_\mu \;\sim\; \tau_{\mathrm{rel}} \times \exp\!\left(\frac{\Delta F}{T}\right),
\label{eq:kramers_time}
\end{equation}
where $\tau_{\mathrm{rel}}$ is the relaxation time of the perpendicular diffusion.  At $\alpha = 0.20$ ($\varphi_{1,\max} \approx 0.57$, $\phi_{\mathrm{eq}} = 0.90$): $v_{\mathrm{entry}} = 0.236$, $\Delta F/T \approx 8.3$ at $N = 50$, giving $\tau \sim 10^5$ MC steps per pattern.

Patterns with lower $\varphi_{1\mu}$ have larger $v_{\mathrm{entry}}$ and exponentially higher barriers.  The escape is dominated by the pattern with the largest overlap, $\varphi_{1\mu} \approx \varphi_{1,\max}$.  With $K$ escape-enabling patterns competing independently, the effective escape time is $\tau_{\mathrm{eff}} \approx \tau_1/K$.

\subsection{Two-pattern thermal equilibrium}
\label{ssec:retrieval}

The geometric centroid~\eqref{eq:centroid} is the $T = 0$ energy minimum.  At finite temperature, the system thermalizes in the two-pattern basin, balancing the energy $u_2(\phi) = -b^{-1}\ln[2(1-b+b\phi)]$ against the entropy of the perpendicular subspace, $s_2(\phi, \varphi_{1\mu}) = \tfrac{1}{2}\ln(1 - 2\phi^2/(1+\varphi_{1\mu}))$.  Setting $\partial(u_2 - T\,s_2)/\partial\phi = 0$ gives the two-pattern self-consistency equation:
\begin{equation}
\frac{1}{1-b+b\,\phi} = \frac{2\,T\,\phi}{(1+\varphi_{1\mu}) - 2\phi^2}\,.
\label{eq:two_pat_sc}
\end{equation}
The equilibrium centroid overlap $\varphi_{\mathrm{cen}}(\varphi_{1\mu}, T)$ depends on both the inter-pattern overlap and the temperature.

The system escapes to the pattern with the largest overlap $\varphi_{1\mu} \approx \varphi_{1,\max}$ (smallest barrier).  Since $\varphi_{1,\max}$ varies slowly with $\alpha$ ($0.51$--$0.67$ for $\alpha = 0.15$--$0.30$) and $\varphi_{\mathrm{cen}}$ is a slowly varying function of $\varphi_{1\mu}$, the centroid peak position is approximately $\alpha$-independent.  At $T = 0.325$: $\varphi_{\mathrm{cen}} \approx 0.82$--$0.84$ across this range, consistent with the observed peak at $\varphi \approx 0.83$ (Fig.~\ref{fig:cdf}).

\subsection{Centroid hopping}
\label{ssec:hopping_theory}

After the system reaches the centroid $(\bxi^1, \bxi^\mu)$, it can escape again by the same mechanism: pattern $\bxi^\mu$ has its own high-overlap neighbors, and one of them, $\bxi^\nu$, can be recruited via a perpendicular rotation in the remaining $(N\!-\!2)$-dimensional subspace.  This time, it is $\bxi^1$ that is pushed out of support while $\bxi^\nu$ enters.  The result is a new two-pattern centroid $(\bxi^\mu, \bxi^\nu)$ at comparable energy:
\begin{equation}
\underbrace{\bxi^1}_{\phi_1 = 0.90}
\;\xrightarrow{\text{1st hop}}\;
\underbrace{(\bxi^1,\,\bxi^\mu)}_{\phi_1 \approx 0.83}
\;\xrightarrow{\text{2nd hop}}\;
\underbrace{(\bxi^\mu,\,\bxi^\nu)}_{\phi_1 \approx 0.48}
\;\xrightarrow{\text{3rd hop}}\;
\underbrace{(\bxi^\nu,\,\bxi^\rho)}_{\phi_1 \approx 0.1\text{--}0.2}
\;\to\; \cdots
\label{eq:hopping_chain}
\end{equation}

After the second hop, the target is no longer in support ($\phi_1 < \varphi_c$).  Its residual overlap with the active spurious pair is
\begin{equation}
\phi_1 = \frac{\varphi_{1\mu} + \varphi_{1\nu}}{\sqrt{2(1+\varphi_{\mu\nu})}}\,.
\label{eq:residual}
\end{equation}
The overlap $\varphi_{1\nu}$ between the target and the newly recruited pattern $\bxi^\nu$ is not independent of $\varphi_{1\mu}$: since $\bxi^\nu$ was selected for having high overlap with $\bxi^\mu$, decomposing $\bxi^\nu$ along $\bxi^\mu$ and perpendicular gives
\begin{equation}
E[\varphi_{1\nu} \mid \varphi_{\mu\nu}] = \varphi_{1\mu}\,\varphi_{\mu\nu}\,.
\label{eq:conditional}
\end{equation}
For $\varphi_{1\mu} \approx \varphi_{\mu\nu} \approx 0.57$ (the most-overlapping neighbors at $\alpha = 0.20$): $\varphi_{1\nu} \approx 0.33$.  Substituting into~\eqref{eq:residual}: $\phi_1 \approx (0.57 + 0.33)/\sqrt{2 \times 1.57} \approx 0.51$, consistent with the observed peak near $\varphi \approx 0.48$ (Fig.~\ref{fig:cdf}).

The dynamics is a random walk on the pattern overlap graph, where each step swaps one member of the active pair for a neighbor of the other.  The two-pattern centroid energy is approximately constant along this walk; only the target's overlap decreases.

\paragraph{No true paramagnetic state.}
In the LSR model, the energy~\eqref{eq:lsr} diverges ($\to +\infty$) when no pattern is in support, because the argument of the logarithm vanishes.  The system is therefore \emph{always} at some centroid.  This can be verified by measuring the maximum alignment with any non-target pattern, $\phi_{\max}^{\mathrm{other}} = \max_{\mu \neq 1} \phi_\mu$: even trials with $\phi_1 \approx 0$ show $\phi_{\max}^{\mathrm{other}} \approx 0.82$, confirming that the system sits at a two-pattern centroid of distant patterns rather than in a paramagnetic state.

\subsection{Predicted escape times and probe-point selection}
\label{ssec:escape_predictions}

Combining the Kramers escape time~\eqref{eq:kramers_time} with the multi-pattern competition ($K$ patterns racing independently), the effective time for the first hop is
\begin{equation}
\tau_{\mathrm{eff}} = \frac{\tau_{\mathrm{rel}}}{K}\,\exp\!\left(\frac{\Delta F}{T}\right),
\label{eq:tau_eff}
\end{equation}
where $\Delta F/T$ is evaluated at $\varphi_{1\mu} = \varphi_{1,\max}$ (the pattern with the lowest barrier) and $K = n(\varphi_{\min})$ is defined in Section~\ref{ssec:escape}.  The relaxation time $\tau_{\mathrm{rel}} \approx (1-\phi_{\mathrm{eq}}^2)N^2/(2.88\,T^2)$ is the Ornstein--Uhlenbeck time scale for diffusion in the entropic potential.  Since the second hop has a comparable barrier (Section~\ref{ssec:escape}), the time to observe $k$ hops is approximately $k\,\tau_{\mathrm{eff}}$.

Table~\ref{tab:escape_times} evaluates~\eqref{eq:tau_eff} across a range of $\alpha$ and $T$.  The barrier $\Delta F/T$ depends weakly on $T$ (it is set primarily by the geometry through $v_{\mathrm{entry}}$ and $\phi_{\mathrm{eq}}$), but the escape time varies over many orders of magnitude due to the exponential.

\begin{table}[h]
\centering
\caption{Predicted first-hop escape time $\tau_{\mathrm{eff}}$ (in MC steps) and number of hops observable within $2^{18} \approx 2.6 \times 10^5$ steps.  Parameters: $b = 2+\sqrt{2}$, $M = 20\,000$.}
\label{tab:escape_times}
\begin{tabular}{ccrccrcl}
\toprule
$\alpha$ & $N$ & $T$ & $\phi_{\mathrm{eq}}$ & $\Delta F/T$ & $K$ & $\tau_{\mathrm{eff}}$ & hops in $2^{18}$ \\
\midrule
0.18 & 55 & 0.325 & 0.899 & 11.0 & 155 & $7 \times 10^5$ & $< 1$ \\
0.18 & 55 & 0.500 & 0.870 & 9.8 & 474 & $4 \times 10^4$ & $\sim 6$ \\
\midrule
0.20 & 50 & 0.200 & 0.928 & 9.4 & 62 & $6 \times 10^5$ & $< 1$ \\
0.20 & 50 & 0.325 & 0.899 & 8.0 & 232 & $2 \times 10^4$ & $\sim 13$ \\
0.20 & 50 & 0.500 & 0.870 & 7.4 & 648 & $2 \times 10^3$ & $\sim 130$ \\
\midrule
0.22 & 45 & 0.150 & 0.942 & 7.3 & 41 & $1 \times 10^5$ & $\sim 2$ \\
0.22 & 45 & 0.325 & 0.899 & 5.8 & 285 & $1.5 \times 10^3$ & $\sim 175$ \\
0.22 & 45 & 0.500 & 0.870 & 5.5 & 728 & $2 \times 10^2$ & $> 10^3$ \\
\midrule
0.24 & 41 & 0.100 & 0.958 & 5.9 & 22 & $8 \times 10^4$ & $\sim 3$ \\
0.24 & 41 & 0.200 & 0.928 & 4.6 & 114 & $2 \times 10^3$ & $\sim 130$ \\
0.24 & 41 & 0.500 & 0.870 & 4.2 & 817 & $44$ & $> 10^3$ \\
\bottomrule
\end{tabular}
\end{table}

Three regimes are apparent:
\begin{enumerate}
\item \textbf{Frozen} ($\tau_{\mathrm{eff}} \gg 2^{18}$): the system remains in the retrieval basin for the duration of the simulation.  This occurs at low $T$ and low $\alpha$ (e.g.\ $\alpha = 0.18$, $T = 0.325$).
\item \textbf{Resolvable} ($\tau_{\mathrm{eff}} \sim 10^3$--$10^5$): individual hops can be observed and timed.  The first-hop time distribution should be approximately exponential (Kramers process).  This is the regime where the theory can be tested quantitatively.
\item \textbf{Instantaneous} ($\tau_{\mathrm{eff}} \ll 10^2$): the system hops many times during equilibration and the trajectory shows free diffusion on the pattern overlap graph.
\end{enumerate}

To validate the Kramers prediction~\eqref{eq:tau_eff}, we select probe points in the resolvable regime where $\tau_{\mathrm{eff}} \sim 10^3$--$10^5$ steps (boldface rows in Table~\ref{tab:escape_times}).  At each point, long MC trajectories ($2^{18}$ steps) with $\varphi(t)$ recorded every 64 steps allow direct measurement of the first-hop, second-hop, and subsequent hop times.  The predicted approximate equality $\tau_{\mathrm{1st\,hop}} \approx \tau_{\mathrm{2nd\,hop}}$ (comparable Kramers barriers) can also be tested.

%%% ─────────────────────────────────────────────────────────────────
\section{Monte Carlo simulations}
\label{sec:mc}

\subsection{Simulation protocol}
\label{ssec:protocol}

We simulate the Boltzmann distribution $\propto \exp(-E(\bx)/T)$ for the LSR energy~\eqref{eq:lsr} using Metropolis--Hastings Monte Carlo on $S^{N-1}(\!\sqrt{N})$.  At each step, a candidate state is proposed as $\bx' = \bx + \sigma\boldsymbol{\eta}$, $\boldsymbol{\eta} \sim \mathcal{N}(0, I_N)$, projected onto the sphere ($\bx' \leftarrow \sqrt{N}\,\bx'/\|\bx'\|$), and accepted with probability $\min(1, e^{-\Delta E/T})$.  The step size $\sigma = 2.4\,T/\sqrt{N}$ is adapted to the temperature and dimension.  Proposals falling outside the support of all patterns ($E = +\infty$) are automatically rejected.

Two types of simulations are reported.

\paragraph{Bulk survey (v8m).}  A grid of $\alpha$ and $T$ values is swept.  At each grid point, $M$ patterns are drawn uniformly on $S^{N-1}(\!\sqrt{N})$, with $N = \lfloor\ln M/\alpha\rfloor$.  For each of $n_{\mathrm{dis}}$ independent disorder samples, two replicas (A, B) are initialized near the target $\bxi^1$ with $\phi_{\mathrm{init}} \sim \mathrm{U}(0.75, 1.0)$ and evolved independently.  Each replica undergoes $N_{\mathrm{eq}} = 2^{15}$ equilibration steps followed by $N_{\mathrm{samp}} = 2^{13}$ sampling steps.  The per-trial observables---time-averaged $\varphi_A$, $\varphi_B$, $q_{\mathrm{EA}} = (1/N_{\mathrm{samp}})\sum_t \bx_A(t)\!\cdot\!\bx_B(t)/N$, and $\phi_{\max}^{\mathrm{other}} = (1/N_{\mathrm{samp}})\sum_t \max_{\mu\neq 1}\phi_\mu(\bx_A(t))$---are recorded per disorder sample.  The data combine two runs: Run~1 ($\alpha \in [0.200, 0.350]$, $M = 20$k--$500$k, $n_{\mathrm{dis}} = 256$--$32$) and Run~2 ($\alpha \in [0.005, 0.550]$, $M = 20$k--$500$k, $n_{\mathrm{dis}} = 256$--$32$).

\paragraph{Targeted escape-time measurements (v10).}  At selected $(\alpha, T)$ points, we record the full trajectory $\varphi(t)$ at each MC step, rather than only the time-averaged overlap.  This allows direct measurement of the first-hop time (retrieval $\to$ centroid), second-hop time (centroid $\to$ spurious centroid), and subsequent hops.  The protocol uses longer runs ($N_{\mathrm{eq}}$ up to $2^{18}$) at fewer grid points, with $n_{\mathrm{dis}} = 256$ disorder samples.  The simulation parameters are chosen to resolve the Kramers escape time $\tau \sim \tau_{\mathrm{rel}}\exp(\Delta F/T)$ predicted in Section~\ref{ssec:escape}.

\subsection{Phase diagram}
\label{ssec:phase_diagram}

Figure~\ref{fig:phase_diagram} shows the disorder-averaged alignment $\avg{\varphi}$ in the $(\alpha, T)$ plane from the v8m data.  The ``retrieval'' region ($\avg{\varphi} > 0.7$, blue) and ``non-retrieval'' region ($\avg{\varphi} \approx 0$, red) are separated by a boundary near $\alpha_{\mathrm{th}} \approx 0.25$.

The theory of Section~\ref{sec:theory} reinterprets this diagram:
\begin{itemize}
\item The ``retrieval'' region below $\alpha_{\mathrm{th}}$ is \emph{metastable}: the system has not yet escaped within the allocated MC time ($N_{\mathrm{eq}} + N_{\mathrm{samp}} \approx 4 \times 10^4$ steps).  With longer runs, $\avg{\varphi}$ would decrease as trials escape to centroids.
\item The ``non-retrieval'' region is \emph{not paramagnetic}: $\phi_{\max}^{\mathrm{other}} \approx 0.82$ throughout, confirming that the system sits at a two-pattern centroid of spurious patterns (Section~\ref{ssec:hopping_theory}).  A true paramagnetic state ($\phi_{\max}^{\mathrm{other}} \to 0$) does not exist in LSR because the energy diverges when no pattern is in support.
\item The boundary at $\alpha_{\mathrm{th}}$ is a \emph{kinetic} crossover, not a thermodynamic phase transition: it marks where the Kramers escape time $\tau$ becomes comparable to the MC time, not where a free-energy minimum disappears.
\end{itemize}

\begin{figure}[t]
\centering
%\includegraphics[width=0.5\textwidth]{panels_v8m/phase_diagram.pdf}
\caption{Disorder-averaged alignment $\avg{\varphi}(\alpha, T)$ from v8m.  Blue: retrieval ($\avg{\varphi} > 0.7$); red: non-retrieval ($\avg{\varphi} \approx 0$).  The apparent boundary near $\alpha \approx 0.25$ is a kinetic crossover, not a thermodynamic phase transition.  [Figure to be generated.]}
\label{fig:phase_diagram}
\end{figure}

\subsection{Quantitative tests of the centroid mechanism}
\label{ssec:tests}

\paragraph{Centroid peak position.}
The two-pattern self-consistency equation~\eqref{eq:two_pat_sc} predicts the centroid overlap $\varphi_{\mathrm{cen}}(\varphi_{1,\max}, T)$ as a function of $\alpha$ (through $\varphi_{1,\max}$) and $T$.  Figure~\ref{fig:centroid_position} compares this prediction with the observed centroid peak position extracted from the per-trial overlap distributions at several $(\alpha, T)$ points.  The weak $\alpha$-dependence (less than $\pm 2\%$ across $\alpha = 0.15$--$0.30$) and the $T$-dependence are both consistent with the theory.

\begin{figure}[t]
\centering
%\includegraphics[width=0.5\textwidth]{panels_v8m/centroid_peak_vs_T.pdf}
\caption{Centroid peak position vs $T$ at several $\alpha$ values.  Solid curves: prediction from Eq.~\eqref{eq:two_pat_sc}.  Points: MC data.  [Figure to be generated.]}
\label{fig:centroid_position}
\end{figure}

\paragraph{$\phi_{\max}^{\mathrm{other}}$ across all populations.}
The centroid hopping mechanism (Section~\ref{ssec:hopping_theory}) predicts $\phi_{\max}^{\mathrm{other}} \approx \varphi_{\mathrm{cen}} \approx 0.82$ for \emph{every} non-retrieval trial, regardless of the target overlap $\varphi$.  Table~\ref{tab:phimax} confirms this: for trials binned by $\varphi$, the average $\phi_{\max}^{\mathrm{other}}$ is $\approx 0.82$ in every bin except the retrieval peak.

\begin{table}[h]
\centering
\caption{Average $\phi_{\max}^{\mathrm{other}}$ for trials in each overlap bin, from v8m data at $\alpha = 0.22$, $T = 0.325$.}
\label{tab:phimax}
\begin{tabular}{lcc}
\toprule
Trial population ($\varphi$ range) & $\avg{\phi_{\max}^{\mathrm{other}}}$ & Interpretation \\
\midrule
Retrieval ($0.85$--$0.95$) & $0.60$ & target only in support \\
Centroid ($0.78$--$0.85$) & $0.82$ & target + spurious centroid \\
Spurious centroid ($0.35$--$0.60$) & $0.82$ & spurious pair, target out \\
Distant centroid ($< 0.20$) & $0.82$ & distant pair, target far \\
\bottomrule
\end{tabular}
\end{table}

\paragraph{Residual overlap at the spurious centroid.}
For trials in the $\varphi \approx 0.48$ population, the observed overlap is compared with the prediction $\phi_1 = (\varphi_{1\mu} + \varphi_{1\mu}\varphi_{\mu\nu})/\sqrt{2(1+\varphi_{\mu\nu})}$ from Eqs.~\eqref{eq:residual}--\eqref{eq:conditional}.  At $\alpha = 0.20$: the predicted value is $0.51$; the observed mean is $0.48$.  The $\sim 6\%$ discrepancy is consistent with the variance of $\varphi_{1\nu}$ around its conditional mean.

\subsection{Escape-time measurements}
\label{ssec:escape_times}

The Kramers barrier~\eqref{eq:kramers_barrier} predicts the escape time from the retrieval basin as a function of $\alpha$, $T$, and $N$.  To test this prediction, we record the full trajectory $\varphi(t)$ at selected $(\alpha, T)$ points using long MC runs.

Figure~\ref{fig:escape_times} shows the distribution of first-hop times (defined as the first time $\varphi$ drops below $0.86$) across disorder samples.  The distribution is approximately exponential, consistent with a Kramers process.  The median escape time is compared with the theoretical prediction $\tau_{\mathrm{eff}}$ from Eq.~\eqref{eq:tau_eff}.

\begin{figure}[t]
\centering
%\includegraphics[width=0.5\textwidth]{panels_v10/escape_time_distribution.pdf}
\caption{Distribution of first-hop times at $\alpha = 0.22$, $T = 0.325$.  Solid line: exponential fit.  Dashed line: Kramers prediction.  [Figure to be generated from targeted simulations.]}
\label{fig:escape_times}
\end{figure}

The v8m data, with its fixed MC time, provide an indirect test: the fraction of trials that have escaped by $N_{\mathrm{eq}} + N_{\mathrm{samp}} \approx 4 \times 10^4$ steps equals $P(\mathrm{escape}) = 1 - \exp(-T_{\mathrm{total}}/\tau_{\mathrm{eff}})$.  At $\alpha = 0.20$, $T = 0.325$: the theory predicts $P \approx 7\%$--$15\%$, compared with the observed centroid fraction of $\sim 7\%$--$10\%$.

The second-hop time (centroid $\to$ spurious centroid) is predicted to be comparable to the first-hop time (Section~\ref{ssec:escape}, comparable barrier heights).  This can be tested by identifying trials that have reached the centroid peak and measuring how long they remain there before hopping further.

\paragraph{Metastability is practically irrelevant.}
The escape time $\tau \sim \exp(cN)$ grows exponentially with the system dimension.  At $N = 50$ and $\alpha = 0.20$: $\tau \sim 10^4$--$10^5$ MC steps.  But $\tau$ increases rapidly with $N$: at $N = 100$, the barrier $\Delta F/T$ roughly doubles, giving $\tau \sim 10^{8}$--$10^{10}$.  For any practical retrieval application with $N \gg 50$, the escape time far exceeds any realistic query time.  The retrieval state, though formally metastable, is effectively stable.

%%% ─────────────────────────────────────────────────────────────────
\section{Discussion and summary}
\label{sec:discussion}

[Three regimes of retrieval stability.  Practical implications: $\tau \sim e^{cN}$.  Comparison with classical Hopfield mixture states.  Open questions.]

\bibliography{references}

\end{document}
